% Generated draft: requires mathematical and source review.
\documentclass[11pt]{article}
\usepackage[margin=25mm]{geometry}
\usepackage{fontspec}
\setmainfont{Latin Modern Roman}
\usepackage{amsmath,amssymb,mathtools,unicode-math}
\setmathfont{Latin Modern Math}
\usepackage{enumitem,needspace,xurl,hyperref}
\setlength{\parindent}{0pt}
\setlength{\parskip}{5pt}
\setlength{\emergencystretch}{3em}
\newcommand{\N}{\mathbb{N}}
\newcommand{\Z}{\mathbb{Z}}
\newcommand{\Q}{\mathbb{Q}}
\newcommand{\R}{\mathbb{R}}
\newcommand{\C}{\mathbb{C}}
\newcommand{\F}{\mathbb{F}}
\newcommand{\A}{\mathbb{A}}
\newcommand{\PP}{\mathbb{P}}
\newcommand{\E}{\mathbb{E}}
\newcommand{\eps}{\varepsilon}
\newcommand{\dd}{\,\mathrm{d}}
\begin{document}
\section*{UnsolvedMath ID: 30005475}
\textbf{Source problem code:} catalogue entry
\subsection{Definitions and mathematical statement}
Under the standard real-matrix interpretation, let \(m,n\geq 1\) be integers and identify \(\R^{m\times n}\) with \(\R^{mn}\). A function \(f:\R^{m\times n}\to\R\) is rank-one convex if, for every \(A,B\in\R^{m\times n}\) with \(\operatorname{rank}(A-B)\leq 1\) and every \(t\in[0,1]\),
\[
f((1-t)A+tB)\leq (1-t)f(A)+tf(B).
\]
For a nonempty finite set \(K\subset\R^{m\times n}\), use the functional definition of its rank-one convex hull:
\[
K^{\mathrm{rc}}=\left\{X\in\R^{m\times n}: f(X)\leq\max_{A\in K}f(A)\text{ for every real-valued rank-one convex }f:\R^{m\times n}\to\R\right\}.
\]
A subset of \(\R^{mn}\) is semialgebraic if it is a finite Boolean combination of sets defined by polynomial equalities and inequalities with real coefficients.

\textbf{Conjecture 2 (standard real-matrix formulation).} For every pair of integers \(m,n\geq1\) and every nonempty finite set \(K\subset\R^{m\times n}\), the set \(K^{\mathrm{rc}}\) is semialgebraic. If the empty finite set is included, adopt the convention \(\varnothing^{\mathrm{rc}}=\varnothing\), which is semialgebraic.

\textbf{Unresolved formulation qualifications.} The supplied excerpt does not explicitly specify the scalar field, matrix sizes, or definition of rank-one convex hull. The real field, universal quantification over rectangular matrix sizes, and functional hull above make a conventional interpretation precise, but require confirmation against the full source. In particular, the smallest set containing \(K\) and closed under segments in rank-one directions, often called a lamination hull, must not be silently identified with the functional hull used here.

\textbf{Supplied context and qualifications.} The excerpt states that the conjecture holds in \(\R^2\) identified with diagonal \(2\times2\) matrices. Its reference [4, Theorem 1.1] is said to establish a directional analogue for cones in \(\R^2\) with finitely many rays and to provide an algorithm for the corresponding hull. Reference [11] is said to give an analogous result when the direction cone consists of \(d\) linearly independent vectors in \(\R^d\); the precise meaning of this cone description is not supplied. Reference [12] is said to provide a different higher-dimensional example of a semialgebraic rank-one convex hull. These are supplied contextual claims with incomplete bibliographic identification.

The excerpt suggests examining the arrangement of cones placed at the points of \(K\) in higher dimensions, where the cones need not be linear and may be unions of algebraic varieties. Its proposed connection between an affirmative answer and an algorithm for computing hulls is explicitly heuristic; a negative answer is described as indicating computational complexity. Neither observation is an additional mathematical conclusion. The excerpt also reports that Saugata Basu pointed out, after the talk, a related lamination-hull construction in cryptography concerning the round complexity of randomized functions, citing [1].

The supplied dated assessment labels the problem partially solved and reports semialgebraicity for diagonal cases, finite-ray directional settings, and finite sets of \(2\times2\) triangular matrices, while describing the universal higher-dimensional conjecture as open. These are supplied literature-status claims, not independently verified results in this draft.
\subsection{Short English statement}
The conjecture asks whether the rank-one convex hull of every finite collection of real matrices can be described by finitely many polynomial equalities and inequalities combined using Boolean operations. The displayed statement uses the functional definition of the hull and quantifies over all matrix sizes; these interpretations need source confirmation. The supplied record reports special cases and labels the general problem partially solved, with its literature assessment dated 2026-08-22.
\subsection{Sources}
\begin{itemize}
\item Local UnsolvedMath registry, permanent ID \(30005475\), title \emph{Semialgebraicity of Finite Rank-One Convex Hulls}, source problem code \texttt{catalogue entry}. The supplied record identifies the passage as Conjecture 2, assigns proposed year 2023, and supplies no named proposer. Its status is \texttt{partially\_solved}. Its statement status is \texttt{corrected\_verified}, with a supplied assertion that the statement was reconstructed and checked against the official report; that verification has not been independently repeated here.
\item \emph{New Directions in Real Algebraic Geometry} (2023), Oberwolfach Reports, pp.\ 813--881, DOI 10.4171/owr/2023/15. Supplied report and source link: \url{https://doi.org/10.4171/owr/2023/15}. The record describes Oberwolfach Reports as documenting mathematical workshops and their open-problem discussions.
\item Supplied literature source: Meroni--Raita, ``semialgebraic rank-one convex hulls'' (the supplied descriptive label), \url{https://arxiv.org/abs/2509.07541}. No fuller bibliographic details are supplied.
\item Supplied literature source labeled ``OWR report'': \url{https://doi.org/10.4171/owr/2023/15}.
\item Literature assessment checked at \textbf{2026-08-22}. This exact date records supplied dated curation, not independent verification. The curation reports the diagonal, finite-ray directional, and finite \(2\times2\) triangular cases, and describes the universal higher-dimensional conjecture as open. It expressly requests independent verification of time-sensitive status and cited claims.
\item The excerpt's reference labels [4, Theorem 1.1], [11], and [12] accompany the special cases described above; labels [6], [7], [8], [10], and [13] are offered for further information on rank-one and directional convexity. Label [1] accompanies the reported cryptographic connection pointed out by Saugata Basu. Their bibliography entries are absent from the supplied record, so these labels cannot be resolved here.
\end{itemize}
\end{document}
